% === Chapter 8: Conclusion ===
\section{Conclusion}
\label{sec:t8}

Through seven problem-driven evolutionary threads,
this book has systematically traced the complete development trajectory of LLM/VLM post-training,
from early instruction tuning to contemporary RL-based reasoning training.
This chapter does not repeat the technical details of each thread,
but instead distills cross-thread commonalities from a macroscopic perspective,
identifies recurring meta-patterns,
and offers an outlook on the future evolution of the post-training landscape.

\subsection{Macroscopic Trends}
\label{sec:t8-trends}

Looking back across the evolutionary trajectories of the seven threads, we can identify seven overarching macroscopic trends.

\paragraph{Trend 1: From supervised imitation to reinforcement learning.}
The technical center of gravity in post-training has undergone a systematic shift from SFT to RL.
In the early stage, SFT was the dominant force in post-training ---
FLAN \cite{wei2022flan}, the SFT stage of InstructGPT \cite{ouyang2022instructgpt},
and open-source projects such as Alpaca \cite{taori2023alpaca} all revolved around supervised learning as their core.
\crossref{1}{sec:t1}
The introduction of RLHF (\crossref{2}{sec:t2}) marked post-training's entry into the phase of optimizing for human preferences,
but early RLHF still heavily relied on SFT warm start, with RL serving more as a ``fine-tuning of fine-tuning.''
The true inflection point emerged in the reasoning thread (\crossref{4}{sec:t4}):
OpenAI o1 \cite{openai2024o1} and DeepSeek-R1 \cite{deepseek2025r1}
demonstrated that large-scale RL training can elicit entirely new reasoning capabilities from models ---
capabilities that cannot be obtained through SFT-based imitation of existing data,
but are instead autonomously discovered through RL's exploration-exploitation mechanism.
The emergence of compute-efficient RL algorithms such as GRPO \cite{shao2024grpo}
(\crossref{7}{sec:t7-compute-rl}) further lowered the barrier to RL training,
expanding the RL-first post-training paradigm from a handful of leading laboratories to the broader community.

\noindent
The deeper logic underlying this trend is as follows:
the capability ceiling of SFT is bounded by the quality of human demonstrations in the training data ---
at best, the model can only learn to do what human annotators are capable of doing.
RL, by contrast, allows the model to explore strategy spaces beyond human demonstrations through environmental feedback (reward signals),
with its capability ceiling theoretically bounded only by the quality of the reward signal.
The transition from an ``imitation ceiling'' to ``reward-guided exploration''
represents a fundamental shift in post-training --- from ``teaching models to do what humans already know how to do''
to ``enabling models to discover strategies that humans have not yet mastered.''

\paragraph{Trend 2: From text-only to multimodal.}
Post-training techniques were initially developed entirely within the text modality,
yet nearly every core technique has undergone migration to multimodal settings:
SFT's instruction tuning was extended by LLaVA \cite{liu2023llava} into
visual instruction tuning (\crossref{5}{sec:t5});
RLHF was applied by RLHF-V \cite{sun2023rlhfv} to reduce multimodal hallucination;
DPO was introduced by LLaVA-RLHF \cite{sun2024llava-rlhf} for preference optimization in VLMs.
This cross-modal migration is not a mere transplantation of techniques ---
multimodal settings introduce new dimensions of challenge such as visual grounding
and cross-modal consistency,
driving substantive extensions and improvements of the original techniques.

\paragraph{Trend 3: From human feedback to automated \& verifiable feedback.}
The source of feedback signals has undergone three generations of evolution.
First generation: purely human annotation (InstructGPT's RLHF) ---
high quality but prohibitively expensive and difficult to scale.
Second generation: AI feedback ---
Constitutional AI \cite{bai2022constitutional} enabled models to self-critique,
RLAIF \cite{lee2023rlaif} used AI to generate preference annotations
(\crossref{7}{sec:t7-synthetic}),
and Zephyr \cite{tunstall2023zephyr} achieved a fully zero-human-annotation alignment pipeline.
\crossref{3}{sec:t3}
Third generation: verifiable reward ---
in domains where correctness can be automatically verified, such as mathematics and programming,
rewards are directly based on ground truth (answer correctness, code passing test cases, etc.),
completely bypassing reward models and human annotation.
The success of GRPO \cite{shao2024grpo} and DeepSeek-R1 \cite{deepseek2025r1}
with verifiable reward (\crossref{4}{sec:t4})
raises a fundamental question:
\textbf{Where are the boundaries of verifiable reward?}
Which capabilities can be formalized into verifiable rewards, and which cannot?

\paragraph{Trend 4: From single-turn instruction following to agentic capability.}
The target behavior of models has evolved from ``answering a single question''
to ``completing multi-step complex tasks within an environment.''
Early SFT focused on single-turn instruction following (\crossref{1}{sec:t1}),
which then expanded to multi-turn dialogue (LIMA's ablation study on 30 dialogue examples demonstrated the importance of this direction,
\crossref{7}{sec:t7-data-selection}),
then to tool use (Toolformer \cite{schick2023toolformer}, \crossref{6}{sec:t6}),
and ultimately evolved into full agentic behavior ---
encompassing task planning, tool selection, error recovery, and long-term memory management.
This evolutionary direction poses new data and evaluation challenges for post-training:
agent training data must include trajectories (rather than independent QA pairs),
and evaluation must be conducted in interactive environments (rather than on static benchmarks).

\paragraph{Trend 5: From prohibitive cost to democratized access.}
The efficiency techniques of Thread~7 have expanded post-training from an exclusive domain of large laboratories
to an activity accessible to the entire research community and industry.
LoRA \cite{hu2022lora} reduced the parameter count by 10,000$\times$,
QLoRA \cite{dettmers2023qlora} made 65B models trainable on a single 48GB GPU
(\crossref{7}{sec:t7-foundation}),
Self-Instruct \cite{wang2023selfinstruct} and UltraChat \cite{ding2023ultrachat}
reduced data acquisition costs by several orders of magnitude,
and GRPO \cite{shao2024grpo} eliminated the most resource-intensive component of PPO --- the value model.
The cumulative effect of this ``efficiency stack'' is profound ---
it not only lowered the barrier to entry but also accelerated the iteration speed of the entire field,
creating a positive feedback loop:
lower cost $\rightarrow$ more research participants $\rightarrow$
faster methodological innovation $\rightarrow$ more efficient techniques $\rightarrow$ lower cost.

\paragraph{Trend 6: Thinking models as the default paradigm.}
The most prominent paradigm shift of 2025--2026 is the evolution of reasoning capability
from ``a specialized ability of dedicated models'' to ``a default mode of general-purpose models'' (\crossref{4}{sec:t4-thinking-models}).
The OpenAI o3 series \cite{openai2025o3}, DeepSeek V3.1/V3.2 \cite{deepseek2025v32},
and Qwen3 \cite{yang2025qwen3}
all embed reasoning as a switchable mode within general-purpose models.
This shift has had profound implications for the post-training pipeline:
SFT is no longer the final step but rather the foundation for cold-starting reasoning capabilities (\crossref{1}{sec:t1-2025}),
RLVR has become the core training method for acquiring reasoning capabilities (\crossref{2}{sec:t2-rlvr}),
and thinking mode fusion has become a new component in the post-training pipeline.

\paragraph{Trend 7: Computer use and the agentic frontier.}
In 2024--2025, LLM agents evolved from calling predefined APIs
to directly operating computer graphical user interfaces (\crossref{6}{sec:t6-computer-use}).
Systems such as Claude Computer Use \cite{anthropic2024computeruse}
and OpenAI Operator \cite{openai2025operator} demonstrated
the ability of AI agents to complete complex tasks in real desktop and browser environments.
Meanwhile, standardized protocols such as MCP \cite{anthropic2024mcp} and
the emergence of Agent SDKs are pushing agents from experimental prototypes toward industrial deployment.
This trend poses entirely new training data and evaluation challenges for post-training:
agent training requires GUI interaction trajectories rather than traditional text-only QA pairs,
and evaluation must be conducted through end-to-end testing in real software environments.

\subsection{Cross-Thread Meta-Patterns}
\label{sec:t8-meta}

Beneath the macroscopic trends described above,
the evolution across threads has repeatedly exhibited several structural meta-patterns.
These patterns are not confined to any single thread
but constitute common regularities that pervade the post-training landscape.

\paragraph{Pattern 1: Simplification Through Reformulation.}
Nearly every thread has undergone a process in which ``complex methods are replaced by more elegant reformulations'':
\begin{itemize}[leftmargin=2em]
  \item Thread~2: PPO (4 models + complex training loop)
    $\rightarrow$ DPO (reformulated as a classification loss, 2 models).
  \item Thread~4: PPO + value model + GAE
    $\rightarrow$ GRPO (eliminating the value model, group-relative reward).
  \item Thread~7: Full fine-tuning (updating all parameters)
    $\rightarrow$ LoRA (low-rank decomposition, 0.01\% of parameters).
\end{itemize}
These simplifications were not achieved by sacrificing performance ---
DPO matches or surpasses PPO, GRPO surpasses PPO, and LoRA matches full fine-tuning ---
but rather by discovering redundant computational structures in the original problem
through \textbf{mathematically equivalent transformations or low-rank approximations}.
This pattern hints at a deeper regularity:
many seemingly complex problems in post-training
have an effective dimensionality far lower than their surface complexity.

\paragraph{Pattern 2: Data Quality $\gg$ Data Quantity.}
Evidence across multiple threads consistently supports the conclusion that ``data quality matters far more than quantity'':
\begin{itemize}[leftmargin=2em]
  \item Thread~1/7: LIMA matched the performance of RLHF models using only 1,000 high-quality examples.
  \item Thread~7: AlpaGasus surpassed the full 52K Alpaca dataset using a 9K high-quality subset.
  \item Thread~7: Phi-1 used textbook-quality synthetic data to train a 1.3B model that outperformed
    models 10$\times$ larger.
  \item Thread~7: Muennighoff et al.\ found that perplexity-based quality filtering
    was more effective than simple deduplication.
\end{itemize}
This pattern forms an elegant resonance with LoRA's low-rank finding:
if the ``information content'' required for alignment is inherently small (the Superficial Alignment Hypothesis),
then a small amount of high-quality data naturally suffices ---
adding low-quality data merely introduces noise in a high-dimensional space
without increasing the effective information content.

\paragraph{Pattern 3: The Teacher $\rightarrow$ Self-Improvement Loop.}
Post-training techniques are evolving from ``reliance on an external teacher'' toward ``model self-improvement'':
\begin{itemize}[leftmargin=2em]
  \item Data generation: Self-Instruct (using the model itself to generate training data)
    $\rightarrow$ Humpback (using the model to generate instructions for web text)
    $\rightarrow$ Self-Play (the model plays against itself to generate preference data).
  \item Feedback: Human feedback $\rightarrow$ AI feedback (RLAIF, CAI)
    $\rightarrow$ Self-reward (the model evaluates its own outputs).
  \item Reasoning: SFT imitating CoT $\rightarrow$
    RL autonomously discovering reasoning strategies (DeepSeek-R1's ``aha moment'').
\end{itemize}
This pattern reveals the ultimate direction of post-training:
\textbf{reducing dependence on external resources (human annotations, strong model APIs)
and enabling models to continuously improve through self-interaction and environmental feedback}.
However, this simultaneously introduces new risks ---
self-improvement loops may amplify existing model biases (bias amplification)
or converge to local optima rather than global optima (mode collapse).

\paragraph{Pattern 4: Cross-Thread Technique Migration.}
The cross-thread citation structure of this book reveals a highly active pattern of technique migration:
reward modeling concepts flowed from Thread~2 to Thread~4 (process reward model),
then back to Thread~2 through GRPO;
SFT's instruction tuning flowed from Thread~1 to Thread~5 (LLaVA) and Thread~6 (Toolformer);
Thread~7's LoRA enabled low-cost training across all other threads.
The rise of the RLVR paradigm in 2025 provides a new exemplary case:
the RLVR concept was systematized in Thread~2 (DAPO \cite{yu2025dapo}),
applied at scale in Thread~4 (reasoning training),
and then rapidly migrated to Thread~5 (R1-VL \cite{zhang2025r1vl}, MM-Eureka \cite{meng2025mmeureka}),
forming a clear migration path of T2 $\rightarrow$ T4 $\rightarrow$ T5.
This dense flow of techniques implies that \textbf{post-training is fundamentally a highly interconnected research network,
rather than a collection of independent research directions}.
Any breakthrough in one thread can produce cascading effects in others ---
for example, GRPO was introduced in Thread~7 as an efficiency innovation,
but its greatest impact manifested in Thread~4 (DeepSeek-R1's reasoning training).

\paragraph{Pattern 5: Evaluation Bottleneck.}
A systemic challenge that pervades all threads is the lag in evaluation methodology:
\begin{itemize}[leftmargin=2em]
  \item Thread~2: Reward model overoptimization ---
    proxy reward continues to rise while true reward declines.
  \item Thread~3: Safety benchmarks cannot cover the long tail of jailbreak strategies.
  \item Thread~4: Static reasoning benchmarks cannot assess models' exploratory reasoning capabilities.
  \item Thread~6: Agent evaluation requires interactive environments,
    and existing benchmarks are far less comprehensive than text-only evaluation.
\end{itemize}
The gap between the rapid iteration of training methods and the relative lag of evaluation methods
is one of the most important systemic challenges facing the post-training field.

\subsection{Future Directions}
\label{sec:t8-future}

Based on the trends and meta-patterns described above,
we outline several key future directions for the post-training field.

\paragraph{Direction 1: Unified Post-Training Pipeline.}
Current post-training is typically a multi-stage sequential process (SFT $\rightarrow$ RM $\rightarrow$ RL,
or SFT $\rightarrow$ DPO),
with each stage having independent data, hyperparameters, and optimization objectives.
An important open question is:
\textbf{Can we design a unified single-stage post-training framework
that simultaneously optimizes multiple objectives including instruction following, safety, and reasoning?}
ORPO \cite{hong2024orpo} (merging SFT and DPO into a single stage) and
Self-Play \cite{chen2024selfplay} (merging data generation and preference optimization)
have taken preliminary steps in this direction,
but there remains a long road toward a truly unified framework.

\paragraph{Direction 2: Extending Verifiable Reward --- From a dichotomy to a consistency spectrum.}
The success of GRPO and DeepSeek-R1 is highly dependent on verifiable reward
(mathematical answer correctness, code execution passing tests, etc.).
However, many important capability dimensions in post-training ---
creative writing, open-ended dialogue, value alignment ---
lack natural verifiable rewards.
The traditional approach divides reward signals into ``verifiable'' and ``non-verifiable'' categories,
but this dichotomy overlooks a critical middle ground:
\textbf{consistency verification constitutes a spectrum from external verification to purely internal signals} ---
external consistency (verifiable reward, answer $=$ ground truth)
$\rightarrow$ evaluator consistency (reward model's preference judgment)
$\rightarrow$ internal consistency (self-consistency across multiple model samples).
The most active explorations of the past two years have been concentrated on the internal consistency end of this spectrum.

\textbf{Self-consistency as reward} has an evolutionary chain that reveals the maturation of this direction.
ScPO \cite{prasad2024scpo} was the first to elevate self-consistency from an inference-time voting mechanism
to a training-time preference signal ---
by identifying the most consistent (winner) and least consistent (loser) responses through multiple sampling,
constructing preference pairs for optimization, and approaching supervised performance on GSM8K and MATH.
However, SRT \cite{sheikh2025srt} discovered a critical failure mode:
when a model uses majority voting consistency as its sole reward,
\textbf{consistency collapse} occurs --- the model learns to maximize consistency rather than correctness,
generating increasingly uniform but potentially incorrect answers,
which is essentially a new form of reward hacking.
Co-Rewarding \cite{corewarding2025} mitigated this issue by introducing cross-view supervision:
cross-verification among multiple model perspectives breaks the self-reinforcing cycle of a single model.
CoVo \cite{zhang2025covo} proposed a more refined dual-signal design ---
simultaneously measuring the \textit{consistency} of reasoning paths (the degree of convergence of intermediate states toward the final answer)
and their \textit{volatility} (the degree of deviation toward alternative answers),
matching or surpassing supervised RL performance without any external labels.

\noindent
Beyond the self-consistency path,
two alternative approaches deserve attention.
JEPO \cite{jepo2025} treats chain-of-thought as a latent variable
and completely bypasses the need for verifiable reward by optimizing Jensen's evidence lower bound ---
matching verifiable-reward-based RL on mathematical tasks and significantly surpassing it on open-domain tasks.
RLPR \cite{yu2025rlpr} uses the model's own token probability for the reference answer as a reward,
substituting internal confidence signals for external verifiers.

\noindent
The core open problem in this direction is:
\textbf{A unified reward-free RL framework does not yet exist.}
ScPO/CoVo relies on sampling consistency, JEPO relies on ELBO approximation, and RLPR relies on token probability ---
these methods each have strengths and weaknesses across different task types,
but lack a theory that unifies them within a single mathematical framework.
More importantly, for open-ended creative tasks (such as creative writing and values-oriented dialogue),
even self-consistency signals are not naturally applicable ---
the relationship between ``consistency'' and ``quality'' of creative responses is far more complex than in mathematical reasoning.

\paragraph{Direction 3: Post-Training for Continual Learning.}
Current post-training is typically a one-shot process ---
once a model is trained, it is deployed and never updated.
But the real world continuously generates new knowledge, new demands, and new safety threats.
\textbf{How can we achieve continual post-training
without forgetting previously acquired capabilities?}
The modular nature of LoRA provides a natural framework for this ---
LoRA adapters for different time points and capability dimensions can be independently trained, composed, and swapped ---
but interference between adapters and compositional generalization
remain unsolved challenges.

\paragraph{Direction 4: Theoretical Foundations of Post-Training.}
Despite the enormous practical success of post-training,
its theoretical understanding remains quite limited.
Core open theoretical questions include:
(1) When does LoRA's low-rank assumption hold, and when does it break down?
What is its relationship to the pre-training data distribution and downstream tasks?
(2) Under what conditions are DPO and RLHF equivalent?
Existing equivalence proofs rely on strong assumptions (such as the Bradley-Terry model) ---
how large is the gap when these assumptions are violated in practice?
(3) What is the fundamental limit of reward model overoptimization?
Do alignment methods exist that do not depend on reward model accuracy?
(4) Convergence guarantees for self-improvement loops ---
under what conditions does model self-training converge to better solutions,
and under what conditions does it lead to mode collapse or bias amplification?

\paragraph{Direction 5: Safety-Aware Efficiency.}
Thread~3 (Safety) and Thread~7 (Efficiency) have largely developed independently,
but a potential tension exists between them:
do efficiency techniques (such as synthetic data, distillation, and LoRA)
inadvertently undermine the effectiveness of safety alignment?
Does LoRA's low-rank constraint limit the model's ability to express complex safety boundaries?
Does the distillation process lose the nuanced safety behaviors of the teacher model?
\textbf{Designing post-training methods that maintain efficiency without sacrificing safety}
is a research direction of both practical significance and theoretical depth.

\paragraph{Direction 6: Meta-Capability Training \& Adaptive Reasoning.}
The preceding directions focus on \textit{what} to train (pipeline, reward, safety);
Direction 6 asks a different question:
\textbf{Can post-training cultivate meta-cognitive capabilities
(self-verification, strategy selection, effort allocation),
and allow users to dynamically trade off between reasoning depth and response latency?}

The impact of post-training on model capabilities can be decomposed into three mechanisms:
\textbf{unlocking potential} (activating capabilities that were compressed but poorly organized during pre-training),
\textbf{assembling strategies} (combining basic capabilities into new behavioral patterns), and
\textbf{connecting tools} (teaching the model to use external tools, \crossref{6}{sec:t6}).
This implies that the capability ceiling of post-training is largely determined by pre-training.
However, \textbf{meta-capability} may be an exception ---
it resembles a \textit{computational pattern} rather than specific knowledge,
and RL's reward signal is naturally suited to reinforcing such behavioral patterns of ``when to adopt which strategy.''

\noindent
\textbf{The empirical boundary between ``unlocking'' and ``creating.''}
Yue et al.\ \cite{yue2025rlvr} (NeurIPS 2025) found that:
while models after RLVR training significantly surpass the base model on pass@1,
when the number of samples $k$ is sufficiently large, the base model's pass@$k$ is actually higher ---
\textbf{RLVR does not create new reasoning paths
but learns to more efficiently select the correct paths that already exist in the base model}.
Du et al.\ \cite{du2025reshaping} (COLM 2025 Outstanding Paper)
provided complementary evidence from a mechanistic interpretability perspective:
post-training does not change where knowledge is stored;
its effects can be decomposed into independent adjustments to truthfulness, refusal behavior, and confidence.
\textbf{Meta-capability may be one of the rare exceptions where post-training truly ``creates''
rather than merely ``unlocks''} ---
because computational strategies like ``when to verify'' and ``when to give up''
lack direct training signals in the next-token prediction objective of pre-training.

\noindent
\textbf{Why RL $>$ SFT for meta-capability.}
Chu et al.\ \cite{chu2025sftrl} (ICML 2025) provided the most direct empirical comparison:
RL-trained models generalize across domains on both rule-based and visual variants,
while SFT performs well only within the training distribution ---
\textbf{SFT memorizes, RL generalizes}.
\textbf{SCoRe} \cite{kumar2024score}
was the first to demonstrate that multi-turn RL can effectively train self-correction,
and revealed two mechanistic reasons for SFT's failure:
distribution mismatch (SFT data reflects the error distribution of the data collection policy
rather than the model's own error distribution) and behavior collapse.
``Does Math Reasoning Improve General Capabilities?'' \cite{mathreasoning2025}
further found that RL training preserves the stability of latent-space geometry,
while SFT distorts the representation space, explaining the geometric basis of RL's cross-domain transfer ability.
Li et al.\ \cite{transfer2025} and Sun et al.\ \cite{sun2025grokking}
provided further evidence from the perspectives of transferability and phase transitions, respectively.

\noindent
\textbf{Adaptive reasoning: From ``whether to reason'' to ``what to do when reasoning goes wrong.''}
The proliferation of thinking models (Trend 6) raises a practical question:
how can users dynamically trade off between reasoning depth and response latency?
Kimi k1.5's long2short training \cite{moonshot2025kimi} and
Qwen3's thinking mode fusion \cite{yang2025qwen3}
have already demonstrated the feasibility of adjustable reasoning budgets within a single model.
Chen et al.\ \cite{chen2025underthinking} systematically quantified
the two failure modes of underthinking and overthinking;
Sui et al.\ \cite{sui2025stopoverthinking} systematized efficient reasoning methods into
three categories of strategies: model-based, reasoning output-based, and input prompt-based.
At the post-training level, adaptive compute allocation now spans three tiers:
\textit{whether to reason} --- AdaptThink \cite{zhang2025adaptthink} trains models to autonomously choose between Think and NoThink modes;
\textit{how deeply to reason} --- SelfBudgeter \cite{selfbudgeter2025}
and AutoThink \cite{tu2025autothink} achieve adaptive token allocation
(50--60\% token compression with no accuracy loss);
\textit{what to do when reasoning goes wrong} --- Self-Backtracking \cite{selfbacktracking2025}
trains models to proactively backtrack along erroneous reasoning paths.
\textbf{Meta-R1} \cite{dong2025metar1} and \textbf{MERA} \cite{ha2025mera}
further add an explicit meta-level cognitive layer,
and S$^2$R \cite{ma2025s2r} trains a unified self-verify + self-correct capability.

\noindent
\textbf{Calibration as the highest-leverage meta-capability.}
Among all meta-capabilities, calibration (a model's accurate estimation of its own correctness)
may be the one with the highest leverage ---
it is both a prerequisite for self-correction and the foundation for effort allocation,
as well as a critical dimension of safety (an overconfident wrong answer is more dangerous than admitting uncertainty).
However, AbstentionBench \cite{abstentionbench2025} found that
reasoning fine-tuning reduces models' abstention ability by an average of 24\% ---
\textbf{reasoning training makes models stronger but also more prone to ``not knowing what they don't know''}.
Leng et al.\ \cite{leng2024tamingoverconfidence}
revealed that the reward model in RLHF exhibits a systematic preference for high-confidence responses.
To address this issue,
RLCR \cite{damani2025rlcr} designs RL rewards based on bounded proper scoring rules (such as the Brier score),
making honest confidence reporting the optimal strategy;
``Rewarding Doubt'' \cite{stangel2025rewardingdoubt}
uses a logarithmic scoring rule to simultaneously penalize over-confidence and under-confidence;
MASA \cite{masa2025} extends calibration to a more general notion of meta-awareness
and demonstrates that this capability can transfer across domains.

\noindent
Taken together, the core direction of meta-capability training is
\textbf{integrating calibration, self-correction, and effort allocation
into a unified meta-cognitive RL framework} ---
so that models can not only ``think'' but also accurately assess whether they are thinking correctly,
determine when deeper thinking is needed,
and recognize when they should candidly say ``I'm not sure.''
At the system deployment level, this points toward
reasoning as a metered resource ---
treating ``thinking tokens'' as an independent computational resource at the inference service layer,
with pricing and scheduling that allow different application scenarios to select the optimal reasoning budget
based on the latency-accuracy tradeoff.

\paragraph{Direction 7: Adaptive Learning Loop --- Observation, Adaptation, and Continuous Evolution.}
A fundamental bottleneck in post-training is the \textbf{``observation problem''}:
scalar loss or scalar reward obscures the dynamic changes across multi-dimensional capabilities.
When training reward continues to rise, the model may be severely degrading on certain dimensions,
but these degradations are masked by the averaging effect of aggregate metrics.
ProbeLLM \cite{huang2026probellm}, through hierarchical MCTS,
proactively searches for model failure regions (see \S\ref{sec:t5-self-evolving} for technical details),
revealing that even the strongest models harbor extensive latent weaknesses.
Quantitative studies on the ``alignment tax'' further confirm this issue:
Lin et al.\ \cite{lin2024alignmenttax} found that RLHF causes different capability dimensions to degrade at different rates,
Huang et al.\ \cite{huang2025safetytax} quantified that safety alignment reduces reasoning capability by 7--31\%,
and Du et al.\ \cite{du2025reshaping}
decomposed the effects of post-training into four dimensions from a mechanistic interpretability perspective.
Toward more efficient diagnostics,
ATLAS \cite{li2025atlas} uses IRT-based adaptive testing to reduce the number of evaluation items by 90\%,
LiveBench \cite{white2024livebench} and AutoDetect \cite{cheng2024autodetect}
integrate evaluation into the training loop,
and Wu et al.\ \cite{wu2024featuretracking} use SAEs to track feature-level dynamics.

\noindent
\textbf{From diagnosis to action: Diagnostic-Driven Adaptive Training.}
Once the capability profile has been observed, how can training be adaptively adjusted?
The core principle is the \textit{closed loop}:
diagnose weaknesses $\rightarrow$ generate targeted data $\rightarrow$ train $\rightarrow$ re-diagnose.
DPE \cite{jia2026dpe}'s $p(1-p)$ analysis reveals a key information-theoretic principle:
\textbf{the most informative training signal comes from samples the model ``half-knows''} ---
that is, samples at the model's decision boundary.
Multiple independent works have converged on the same conclusion from different mathematical frameworks:
AERO \cite{gao2026aero}'s entropy-based positioning,
VCRL \cite{jiang2025vcrl}'s reward variance proxy,
DGPO \cite{dai2026dgpo}'s difficulty-aware reweighting,
and Actor-Curator \cite{gu2026actorcurator}'s bandit formulation ---
all methods converge on \textbf{sampling near the model's decision boundary to maximize information gain}.
The ``Difficulty Shift'' phenomenon reported by Zhang et al.\ \cite{zhang2025adcl}
(static difficulty labels rapidly become obsolete as model capabilities evolve)
further argues for the necessity of closed-loop diagnostics.
Beyond the decision boundary, \textbf{high-confidence errors} constitute a complementary high-value training signal:
ACE \cite{ace2026} applies asymmetric penalties to overconfident errors,
forcing the model to re-examine cognitive blind spots where it is ``certain but actually wrong'';
LENS \cite{feng2025lens} extracts training signals from groups of universally failed samples,
enabling the model to learn even from ``total failures.''
DPE's $p(1-p)$ region and ACE/LENS's high-confidence error region
together paint a more complete picture of the training signal distribution:
\textbf{the decision boundary provides maximum information gain,
while high-confidence error regions correct the most dangerous cognitive illusions}.

\noindent
\textbf{Continuous evolution: Breaking through static data saturation.}
Once the training signal from a fixed dataset is exhausted, how can a model continue to evolve?
Naive self-play collapses ---
when a model plays against itself, it tends to converge to trivial equilibria
\cite{liu2026selfplayinfo,sheikh2025srt,yi2025collapse}.
Solutions can be categorized into three types:
corpus-grounded self-play (SPICE \cite{liu2025spice}, SPELL \cite{yang2025spell}),
meta-RL teacher (SOAR \cite{sundaram2026soar}, which generates problems positioned just above the student's capability boundary),
and replay buffer engineering (ExGRPO \cite{zhan2025exgrpo}).

\noindent
The three layers described above --- observation, adaptation, and continuous evolution --- form a complete adaptive learning loop:
ProbeLLM $\approx$ perception system (localizing failure modes),
DPE $\approx$ action system (generating targeted training data),
and self-play mechanisms $\approx$ continuous evolution engine.
However, this loop still has three critical gaps:
(1) Systematic blind spots of the diagnostician itself ---
if the probing model is equally weak on a certain capability dimension, it cannot effectively diagnose that dimension;
(2) Representation of continuous capability spaces ---
predefined categories may miss cross-dimensional weaknesses;
(3) Reward mechanisms for non-verifiable tasks ---
when there is no ground truth, the ``correctness'' of the diagnosis itself is difficult to define.

\paragraph{Direction 8: Multi-Objective Post-Training --- Pareto Optimization.}
Post-training simultaneously optimizes multiple mutually interfering objectives
(helpfulness, safety, reasoning, calibration).
Emerging solutions point toward \textbf{orthogonal subspace decomposition} ---
projecting the gradients of different objectives into mutually non-interfering subspaces.
PAMA \cite{he2025pama} provides an $O(n)$ closed-form Pareto solution,
OrthAlign \cite{orthalign2025} achieves 34--51\% improvement on individual dimensions without harming others,
and NSPO \cite{nspo2025} performs safety optimization within the null space of capability.
The maturation of this direction will transform post-training from
``a tradeoff where improving one objective compromises another''
into ``systematic engineering that precisely navigates the Pareto frontier.''

\bigskip
\noindent
In summary, post-training is evolving from an SFT-centric ``formatting'' stage
into an RL-driven ``capability discovery'' stage.
The underlying logic of this transition is clear:
pre-training endows models with a vast knowledge base and linguistic abilities,
and the task of post-training has deepened from ``teaching models how to present these abilities'' (the SFT paradigm)
to ``enabling models to discover new behavioral strategies through reward-guided exploration'' (the RL paradigm).
The proliferation of thinking models and the breakthroughs in computer use in 2025--2026 further confirm this direction:
RL can not only elicit reasoning capabilities but also train models to execute complex multi-step operations in real environments.
At the same time, continuous advances in efficiency techniques have made these powerful but expensive training paradigms
accessible to a much broader range of research and application.
From the early explorations of RLHF in 2017 to the establishment of thinking models as the default paradigm in 2026,
the continuous cross-pollination and fusion of the seven evolutionary threads suggest that
\textbf{the future of post-training will no longer be a simple combination of independent techniques,
but rather a unified, multi-objective, continuously evolving training system}.
